\subsection{可解圈拓扑群的万有覆叠}

设 G 是局部道路连通的拓扑群。G 的平移都是同胚（TG, III, 第~2~页）。空间 G 是可解圈的条件（IV, 第~340~页，定义 2）等价于 G 具有以下性质：

存在单位元 e 的一个邻域 V，使得由规范嵌入得到的从 \(\pi_{1}(V,e)\) 到 \(\pi_{1}(G,e)\) 的同态的像约化为单位元。

命题 6. --- 设 G 是连通的可解圈拓扑群。存在一个拓扑群 \(\widetilde{\mathbf{G}}\)，其单位元为 \(\widetilde{e}\)，以及满足下列断言的连续同态 \(p\colon \widetilde{\mathbf{G}}\to \mathbf{G}\)：

\begin{enumerate}
\def\labelenumi{\alph{enumi})}
\item
  空间 \(\widetilde{G}\) 是单连通且道路单连通的。赋予其映射 p 后，带基点空间 ( \(\widetilde{G},\widetilde{e}\) ) 是带基点空间 (G,e) 的万有覆叠。
\item
  p 的核 N 是 \(\widetilde{G}\) 的离散子群，包含于 \(\widetilde{G}\) 的中心。同态 \(\mathrm{N}\to\mathrm{Aut}_{\mathrm{G}}(\widetilde{\mathrm{G}})\) 将 \(n\in N\) 送到 \(\widetilde{G}\) 中的右平移，是群同构。从 \(\pi_{1}(\mathrm{G})\) 到 N 的同态把 \(\gamma\in\pi_{1}(\mathrm{G})\) 送到满足 \(\widetilde{e}\cdot\gamma=n\) 的 N 中唯一元素 n，是群同构。
\item
  如果 G' 是以 e' 为单位元的拓扑群，且 p': G' → G 是使 G' 成为 G 的覆叠的群同态，那么唯一连续映射 u: \(\widetilde{G} \rightarrow G'\) 满足 \(u(\widetilde{e}) = e'\) 且 \(p' \circ u = p\)，并且 u 是群同态。赋予 \((\widetilde{G}, \widetilde{e})\) 以映射 u 后，它是 \((G', e')\) 的万有覆叠。
\end{enumerate}

根据 IV, 第~342~页的定理 1，存在 G 的覆叠 \((\widetilde{\mathbf{G}},p)\)，它单连通且道路单连通，是群 \(\pi_1(\mathbf{G})^\circ\) 的伽罗瓦覆叠；另有一点 \(\widetilde{e}\) 属于 \(p^{-1}(e)\)，使得带基点覆叠 \((\widetilde{\mathbf{G}},\widetilde{e})\) 是 \((\mathbf{G},e)\) 的万有覆叠。

根据 IV, 第~375~页的命题 4 和 IV, 第~377~页的命题 5，在空间 \(\widetilde{\mathrm{G}}\) 上存在唯一的与其拓扑相容的群结构，使 \(p\) 成为群同态且 \(\widetilde{e}\) 成为单位元。由命题 4 可知，群 \(\widetilde{\mathrm{G}}\) 满足断言 \(a)\) 和 \(b)\)。

证明断言 \(c\)。设 \(\mathbf{G}'\) 是拓扑群，且 \(p'\colon \mathbf{G}' \to \mathbf{G}\) 是使 \(\mathbf{G}'\) 成为 \(\mathbf{G}\) 的覆叠的群同态。令 \(e'\) 为 \(\mathbf{G}'\) 的单位元。存在唯一的映射 \(u\colon \widetilde{\mathbf{G}} \to \mathbf{G}'\)，满足 \(p = p' \circ u\) 且 \(u(\widetilde{e}) = e'\)，这是因为 \((\widetilde{\mathbf{G}}, \widetilde{e})\) 是带基点空间 \((\mathbf{G}, e)\) 的万有覆叠。由于映射 \(p\) 和 \(p'\) 是满射群同态，\(u\) 是群同态。在映射 \(u\) 下，\(\widetilde{\mathbf{G}}\) 是 \(\mathbf{G}'\) 的覆叠（I, 第~81~页，命题 7），因此 \((\widetilde{\mathbf{G}}, \widetilde{e})\) 是 \((\mathbf{G}', e')\) 的万有覆叠（IV, 第~345~页，定理 1 的推论 2）。

在本命题的假设下，滥用术语地说，\(\tilde{G}\) 是 \(G\) 的万有覆叠。

例。--- 当 G 是 \(\mathbf{R}\) 或 \(\mathbf{C}\) 上的连通李群时，命题 6 尤其适用。此时，在 \(\widetilde{\mathbf{G}}\) 上存在唯一的解析流形结构，使投影 \(p\colon\widetilde{\mathbf{G}}\to\mathbf{G}\) 成为解析流形的平展态射（VAR R, §1, 5.8.2, 第~48~页）。对于此流形结构，\(\widetilde{\mathbf{G}}\) 是李群（LIE, III, 第~113~页，推论）。

附注。--- 设 G 是连通的可解圈拓扑群，e 是其单位元。设 \(\widetilde{G}\) 是道路单连通的拓扑群，单位元为 \(\widetilde{e}\)，且 \(p\colon\widetilde{G}\to G\) 是使 \(\widetilde{G}\) 成为 G 的万有覆叠的群同态。令 N 表示 p 的核。

右平移 \(\delta_{h}\)（分别为左平移）由元素 \(h \in N\) 给出，是主覆叠 \(\widetilde{G}\) 的 G-自同构，且映射 \(h \mapsto \delta_{h}\) 是从群 N 到此主覆叠的自同构群的同构。

对于 N 的每个子群 K，由 p 得到的映射 \(p'\colon\widetilde{G}/K\to G\) 是 G 的伽罗瓦覆叠，且 \(\operatorname{Aut}_{\mathrm{G}}(\widetilde{\mathrm{G}}/\mathrm{K})\) 可识别为群 N/K。当把群 N 与 \(\pi_{1}(\mathrm{G})\) 识别时，群 \(p'_{*}(\pi_{1}(\widetilde{\mathrm{G}}/\mathrm{K}))\) 可识别为 N 的子群 K。此外，\(\widetilde{G}\) 是 \(\widetilde{G}/K\) 的覆叠，因为 G 局部连通（I, 第~81~页，命题 7）。

反之，G 的每个非空连通覆叠 E 都 G-同构于这种类型的覆叠（I, 第~113~页，定理 3 和 I, 第~111~页，命题 10）。事实上，取纤维 \(E_{e}\) 中的一点 x，并令 f 为从 \(\widetilde{G}\) 到 E 的唯一同态，它将 \(\widetilde{e}\) 映到 x。赋予 \(\widetilde{G}\) 以 f 后，它是 E 的伽罗瓦覆叠；子群 \(\mathrm{Aut}_{\mathrm{E}}(\widetilde{\mathrm{G}})\)

是 \(\mathrm{Aut}_{\mathrm{G}}(\widetilde{\mathrm{G}})\) 的子群，可识别为 \(f^{-1}(x)\)，因而是 N 的子群。于是，\(f\) 诱导出从 \(\widetilde{\mathrm{G}} / f^{-1}(x)\) 到 E 的 G-同构。通过结构转移，这在 E 上给出一个拓扑群结构，使 \(x\) 成为单位元，并使映射 \(f\) 成为满射同态。G-空间 E 的投影于是是群同态，因此 E 的合成法则就是 IV, 第~375~页的命题 4 所定义的合成法则。

设 \((\mathrm{E},q)\) 和 \((\mathrm{E}^{\prime},q^{\prime})\) 是 G 的连通覆叠，令 x 为 \(E_{e}\) 中的一点，令 \(x^{\prime}\) 为 \(E_{e}^{\prime}\) 中的一点。存在从 E 到 \(E^{\prime}\) 的 G-态射的充要条件是 \(p_{*}(\pi_{1}(\mathrm{E},x))\) 包含于 \(p_{*}^{\prime}(\pi_{1}(\mathrm{E}^{\prime},x^{\prime}))\)。若此条件满足，则存在唯一的 G-态射 \(f\colon E\to E^{\prime}\)，使得 \(f(x)=x^{\prime}\)（III, 第~311~页，命题 1 的推论 2）。若给 E 和 \(E^{\prime}\) 赋予群合成法则，使 q 和 \(q^{\prime}\) 成为同态且 x 和 \(x^{\prime}\) 成为单位元，则映射 f 是群同态。事实上，f 可识别为规范同态 \(\widetilde{\mathrm{G}}/p_{*}(\pi_{1}(\mathrm{E},x))\to\widetilde{\mathrm{G}}/p_{*}^{\prime}(\pi_{1}(\mathrm{E}^{\prime},x^{\prime}))\)。

命题 7. --- 两个连通的可解圈拓扑群局部同构的充要条件是它们的万有覆叠群作为拓扑群同构。

连通的可解圈拓扑群与其万有覆叠局部同构（IV, 第~369~页，命题 1），所以该条件是充分的。根据 IV, 第~372~页命题 2 的推论 2，该条件是必要的，因为连通的可解圈拓扑群的万有覆叠是单连通的（IV, 第~379~页，命题 6）。

命题 8. --- 设 G 是连通的可解圈拓扑群，且 \(\widetilde{\mathbf{G}}\) 是 G 的万有覆叠。设 V 是 G 的单位元 e 的连通开邻域，使得由规范同态得到的从 \(\pi_1(\mathrm{V} \cdot \mathrm{V}, e)\) 到 \(\pi_1(\mathrm{G}, e)\) 的像约化为单位元。令 F(V, r) 表示由生成集 V 和关系元集合 r（关系元为 \(xyz^{-1}\)）定义的群，其中 \((x, y, z)\) 遍历 V × V × V 中满足 \(xy = z\) 的元素。令 j: V → F(V, r) 为规范映射。

存在唯一的同构 \(f\)，它把群 \(\mathrm{F}(\mathrm{V},\mathbf{r})\) 同构到 \(\widetilde{\mathrm{G}}\)，并且 \(f \circ j\) 是从 V 到 G 的规范嵌入到 \(\widetilde{\mathrm{G}}\) 的提升。

集合 V · V 连通且开，因而局部道路连通。令 p 为 \(\widetilde{G}\) 的投影。存在 p 在 V · V 上的连续截面 s，使得 \(s(e) = \widetilde{e}\)，其中 \(\widetilde{e}\) 表示 \(\widetilde{G}\) 的单位元（III, 第~308~页，命题 1）。置 \(\widetilde{V} = s(V)\)；集合 \(\widetilde{V}\) 是 \(\widetilde{e}\) 在 \(\widetilde{G}\) 中的连通开邻域，且 s 是从 V · V 到 \(\widetilde{V} \cdot \widetilde{V}\) 的同胚。映射 \((x, y) \mapsto s(x)s(y)\) 和 \((x, y) \mapsto s(xy)\) 是到 \(\widetilde{G}\) 的提升，提升映射 \((x, y) \mapsto xy\) 从 V × V 到 G，并在 \((e, e)\) 处重合；由于 V × V 连通，它们在 V × V 上重合（I, 第~34~页，推论 1）。此外，若 \((\widetilde{x}, \widetilde{y}, \widetilde{z}) \in \widetilde{V} \times \widetilde{V} \times \widetilde{V}\)，则条件 \(\widetilde{x}\widetilde{y} = \widetilde{z}\) 和 \(p(\widetilde{x})p(\widetilde{y}) = p(\widetilde{z})\) 等价。因此，同构 \(f: F(V, r) \to \widetilde{G}\) 的存在性由命题 2 的推论 3（IV, 第~372~页）得出。

设 \(g\) 是满足本命题条件的 \(\mathrm{F}(\mathrm{V},\mathbf{r})\) 的同构。等式 \(e \cdot e = e\) 蕴含 \(eee^{-1} \in \mathbf{r}\)，所以 \(j(e)\) 是 \(\mathrm{F}(\mathrm{V},\mathbf{r})\) 的单位元。由于 V 连通，\(g \circ j\) 是从 V 到 G 的规范嵌入到 \(\widetilde{\mathrm{G}}\) 的唯一连续提升。因此，\(f\) 和 \(g\) 在 \(j(\mathrm{V})\) 上重合。由于 \(j(\mathrm{V})\) 生成群 \(\mathrm{F}(\mathrm{V},\mathbf{r})\)，有 \(f = g\)。

